\documentclass[../main.tex]{subfiles}

\begin{document}

\ifSubfilesClassLoaded{

    %\tableofcontents
    
    \setcounter{chapter}{6}
}{}

\chapter{ HW7 }
 代靖涵 25120222201319

 \begin{exercise}{}{}
    设$X$为仅取非负整数的离散随机变量,若其数学期望存在,证明:
\begin{enumerate}
\item $E(X) = \sum_{k=1}^{\infty} P(X \geq k);$
\item $\sum_{k=0}^{\infty} kP(X > k) = \frac{1}{2}(E(X^2) - E(X)).$
\end{enumerate}
 \end{exercise}

 \begin{proof}
    \begin{enumerate}
        \item  \[
        \begin{aligned}
        E\left( X \right)&= \sum _{k= 1}^{\infty}k P\left( X= k \right)  \\ 
         &= \sum _{k= 1}^{\infty}\sum _{j= 1}^{k}P\left( X= k \right)\\ 
          &= \sum _{j= 1}^{\infty}\sum _{k= j}^{\infty}P\left( X= k \right) \\ 
           &= \sum _{j= 1}^{\infty}P\left( X\ge j \right)= \sum _{k= 1}^{\infty}P\left( X\ge k \right)  
        \end{aligned} 
        \] 
        \item  \[
       \begin{aligned}
       \sum _{k= 0}^{\infty}kP\left( X> k \right)&=\sum _{k= 1}^{\infty}\sum _{j= k+ 1}^{\infty}kP\left( X= j \right)\\ 
        &=   \sum _{k= 1}^{\infty}\sum _{j= k}^{\infty}\left( k-1 \right) P\left( X= j \right) \\ 
        &= \sum _{j= 1}^{\infty}\sum _{k= 1}^{j}\left( k-1 \right)P\left( X= j \right)\\ 
         &=   \sum _{j= 1}^{\infty} \frac{1 }{2 }j\left( j-1 \right)P\left( X= j \right)   \\ 
          &= \frac{1}{2}\sum _{j= 1}^{\infty}j^{2}P\left( X= j \right)-\frac{1 }{2 }\sum _{j= 1}^{\infty}jP\left( X= j \right)\\ 
           &= \frac{1}{2}E\left( X^{2} \right)-\frac{1}{2}E\left( X \right)     
       \end{aligned}
        \]
    \end{enumerate}
    
 \end{proof}
 \begin{exercise}{}{10-20-1}
    设连续随机变量$X$的分布函数为$F(x)$,且数学期望存在,证明:
$E(X) = \int_{0}^{\infty} [1 - F(x)] dx - \int_{-\infty}^{0} F(x) dx.$
 \end{exercise}
\begin{proof}
    设 \(  X^{+ }  \)和 \(  X^{-}  \)分别是 \(  X  \)的正部和负部.   注意到 \[
    X^{+ }\left(  \omega  \right)= \int_{0}^{X^{+ }\left(  \omega  \right) }1\,\mathrm{d} x= \int_{0}^{\infty}\chi _{\left\{ X^{+ }\left(  \omega  \right)>  x  \right\} }\left( x \right)\,\mathrm{d} x  
    \]对两边取期望, 我们得到 \[
    E\left( X^{+ } \right)= \int \left( \int_{0}^{\infty}\chi _{\left\{ X^{+ }\left(  \omega  \right)> x  \right\} }\left( x \right)\,\mathrm{d} x  \right)\,\mathrm{d} P  
    \]它等于在乘积空间\(  \mathbb{R} ^{+ }\times  \Omega   \)上 对非负函数 \(  f\left( x, \omega  \right)= \chi _{\left\{ X^{+ }\left(  \omega  \right)> x  \right\}} \)的积分, 由Tonelli定理 \[
   \begin{aligned}
    E\left( X^{+ } \right)&=  \int_{0}^{\infty} \int \chi _{\left\{ X^{+ }\left(  \omega  \right)>  x  \right\} }\,\mathrm{d} P\,\mathrm{d} x   \\ 
     &= \int_{0}^{\infty} E\left(\chi _{\left\{ X^{+ }>  x \right\}}  \right)\,\mathrm{d} x\\ 
      &= \int_{0}^{\infty} E\left( 1-\chi _{\left\{ X^{+ }\le x \right\}} \right)\,\mathrm{d} x\\ 
       &= \int_{0}^{\infty}1- E\left( \chi _{\left\{ X^{+ }\le x \right\}} \right)\,\mathrm{d} x\\ 
        &= \int_{0}^{\infty}\left( 1- E\left( \chi _{\left\{ X\le x \right\}} \right)  \right)\,\mathrm{d} x    \\ 
         &= \int_{0}^{\infty}\left( 1-F\left( x \right)  \right)\,\mathrm{d} x 
   \end{aligned}
    \] 另一方面 \[
    X^{-}\left(  \omega  \right)= \int_{0}^{X^{-}\left(  \omega  \right) }1\,\mathrm{d} x= \int_{0}^{\infty}\chi_{ \left\{ X^{-}\left(  \omega  \right)> x  \right\}}\,\mathrm{d} x  
    \]对两侧取期望, 得到 \[
    E\left( X^{-} \right)= \int \left( \int_{0}^{\infty}\chi _{\left\{ X^{-}\left(  \omega  \right)> x  \right\}}\,\mathrm{d} x \right)  \,\mathrm{d} P
    \]由Tonelli定理,  \[
    \begin{aligned}
    E\left( X^{-} \right)&=  \int_{0}^{\infty}  \int \chi _{\left\{ X^{-}\left(  \omega  \right)> x  \right\}}\,\mathrm{d} P   \,\mathrm{d} x\\ 
     &= \int_{0}^{\infty}E\left( \chi _{\left\{ X^{-}\left(  \omega  \right)  > x\right\}} \right)\,\mathrm{d} x\\ 
      &= \int_{0}^{\infty}E\left( \chi _{\left\{ X\le -x \right\}} \right)\,\mathrm{d} x\\ 
       &= \int_{0}^{\infty}F\left( -x \right)\,\mathrm{d} x\\ 
        &= \int_{-\infty}^{0}F\left( x \right)\,\mathrm{d} x     
    \end{aligned}
    \]最后, 由期望的线性,  \[
    \begin{aligned}
    E\left( X \right)&= E\left( X^{+ }-X^{-} \right)\\ 
     &= E\left( X^{+ } \right)-E\left( X^{-} \right)\\ 
      &= \int_{0}^{\infty}\left( 1-F\left( x \right)  \right)\,\mathrm{d} x-\int_{-\infty}^{0}F\left( x \right)\,\mathrm{d} x       
    \end{aligned}
    \]
\end{proof}
 \begin{exercise}{}{10-21-1}
    设$X$为非负连续随机变量,若$E(X^n)$存在,试证明:
\begin{enumerate}
\item $E(X) = \int_{0}^{\infty} P(X > x) dx;$
\item $E(X^n) = \int_{0}^{\infty} nx^{n-1} P(X > x) dx.$
\end{enumerate}
 \end{exercise}
\begin{proof}
    \begin{enumerate}
        \item 由Lyapunov's 不等式,  \[
        E\left( X \right)\le \left( E\left( X^{n} \right)  \right)^{\frac{1 }{n } }< \infty  
        \]因此期望存在. 由Exercise \ref{ex:10-20-1}, \[
        E\left( X \right)= \int_{0}^{\infty}\left( 1-F\left( x \right)  \right)\,\mathrm{d} x-\int_{-\infty}^{0}F\left( x \right)\,\mathrm{d} x   
        \]由于 \(  X  \)非负, 因此 \[
        F\left( x \right)= P\left( X\le x \right)= 0, \forall x\le 0  
        \] 从而 \[
        \int_{-\infty}^{0}F\left( x \right)\,\mathrm{d} x= \int_{-\infty}^{0}0\,\mathrm{d} x= 0 
        \] 另一方面,  \[
        1-F\left( x \right)= 1-P\left( X\le x \right)= P\left( X> x \right)   
        \]故 \[
        E\left( X \right)=  \int_{0}^{\infty}P\left( X> x \right)\,\mathrm{d} x  
        \]
        \item  注意到  \[
      X^{n}\left(  \omega  \right)= \int_{0}^{X\left(  \omega  \right) }n x^{n-1}\,\mathrm{d} x= \int_{0}^{\infty} \chi _{\left\{ X\left(  \omega  \right)> x  \right\}}n x^{n-1}\,\mathrm{d} x
        \]对两边取期望, 得到 \[
        E\left( X^{n} \right)= \int \left( \int_{0}^{\infty} {\chi _{\left\{ X> x \right\}}} n x^{n-1}\,\mathrm{d}x \right)\,\mathrm{d} P  
        \]由Tonelli定理, \[
        \begin{aligned}
        E\left( X^{n} \right)&= \int_{0}^{\infty} \left( \int \chi _{\left\{ X> x \right\}} nx^{n-1}\,\mathrm{d} P\right)\,\mathrm{d} x\\ 
         & = \int_{0}^{\infty} nx^{n-1} \int  \chi _{\left\{ X> x \right\}}\,\mathrm{d} P\,\mathrm{d} x \\ 
          &=  \int _{0}^{\infty}nx^{n-1}P\left( X> x \right)\,\mathrm{d} x\\ 
           &=  \int_{0}^{\infty} nx^{n-1}P\left( X> x \right)\,\mathrm{d} x  
        \end{aligned}  
        \]
    \end{enumerate}
    
\end{proof}
 \begin{exercise}{}{}
设随机变量$X$满足$E(X) = \text{Var}(X)=\lambda$,已知$E[(X-1) (X-2)] = 1$,试求$\lambda$.
\end{exercise}
\begin{proof}
     \[
     \lambda =  \mathrm{Var}\left( X \right)= E\left( X^{2} \right)-\left( E\left( X \right)  \right)^{2}= E\left( X^{2} \right)- \lambda ^{2}    
     \]因此 \[
     E\left( X^{2} \right)=  \lambda ^{2}+  \lambda  
     \]由期望的线性性,  \[
   \begin{aligned}
     1= E\left[ \left( X-1 \right)\left( X-2 \right)   \right] &= E\left( \left( X^{2}-3X+ 2 \right)  \right)\\ 
      &= E\left( X^{2} \right)-3E\left( X \right)+ 2\\ 
       &=  \lambda ^{2}+  \lambda -3 \lambda + 2\\ 
        &=     \left(  \lambda -1 \right)^{2}+ 1 
   \end{aligned}
     \]因此 \[
     \left(  \lambda -1 \right)^{2}= 0\implies  \lambda = 1 
     \]
\end{proof}
\begin{exercise}{}{}
假设有10只同种电器元件,其中有两只不合格品,装配仪器时,从这批元件中任取一只,如是不合格品,则扔掉重新任取一只;如仍是不合格品,则扔掉再取一只,试求在取到合格品之前,已取出的不合格品数的方差.
\end{exercise}
\begin{proof}
    设取出不合格品的数量为 \(  X  \). 则 \[
    P\left( X= 0\right)=  \frac{4 }{5 },\quad P\left( X= 1 \right)= \frac{1 }{5 }\cdot \frac{8 }{9 },\quad P\left( X= 2 \right)= \frac{1 }{5 }\cdot \frac{1 }{9 }       
    \] 我们有 \[
    E\left( X \right)= 0\cdot P\left( X= 0 \right)+  P\left( X= 1 \right)+ 2P\left( X= 2 \right)   = \frac{2 }{9 } 
    \] \[
    E\left( X^{2} \right)= 0\cdot P\left( X= 0 \right)+ 1\cdot P\left( X= 1 \right)+ 4P\left( X= 2 \right)= \frac{4 }{15 }     
    \]于是 \[
    \mathrm{Var}\left( X \right)= E\left( X^{2} \right)- \left( E\left( X \right)  \right)^{2}=   = \frac{88 }{405 } 
    \]
\end{proof}
\begin{exercise}{}{}
已知 $E(X)=-2$, $E(X^2) = 5$,求 $\text{Var}(1-3X)$.
\end{exercise}
\begin{proof}
    \[
    \mathrm{Var}\left( X \right)= E\left( X^{2} \right)-\left( E\left( X \right)  \right)^{2}= 1   
    \] \[
    \mathrm{Var}\left( 1-3X \right)= \left( -3 \right)^{2}\mathrm{Var}\left( X \right)= 9   
    \]
\end{proof}
\begin{exercise}{}{}
设$P(X=0)=1-P(X=1)$,如果$E(X)=3\text{Var}(X)$,求$P(X=0)$.
\end{exercise}
\begin{proof}
     我们有 \[
     P\left( X= 0 \right)+ P\left( X= 1 \right)= 1  
     \]这表明 \[
     \left\{ X\neq 0 \right\}\cup \left\{ X\neq 1 \right\}
     \]是一个零测集. 因此 \[
     E\left( X \right)=  \int X\,\mathrm{d} P=  \int X\left( \chi _{\left\{ X= 0 \right\}} + \chi _{\left\{ X= 1 \right\}}\right)\,\mathrm{d} P= \int \chi _{\left\{ X= 1 \right\}} \,\mathrm{d} P= P\left( X= 1 \right) 
     \]此外 \[
     E\left( X^{2} \right)=  \int X ^{2}\left( \chi _{\left\{ X= 0 \right\}}+ \chi _{\left\{ X= 1 \right\}} \right)\,\mathrm{d} P= \int \chi _{\left\{ X= 1 \right\}}\,\mathrm{d} P= P\left( X= 1 \right)   
     \]因此\[
    P\left( X= 1 \right)=   E\left( X \right)=  3 E\left( X^{2} \right)- 3 \left( E\left( X \right)  \right)^{2}   =  3P\left( X= 1 \right)- 3 \left( P\left( X= 1 \right)  \right)^{2}  
     \]解得 \[
     P\left( X= 1 \right)= \frac{2 }{3 }  ,\quad \text{或者},\quad P\left( X= 1 \right)= 0 
     \]从而 \[
     P\left( X= 0 \right)= \frac{1 }{3 },\quad \text{或者} ,\quad P\left( X= 0 \right)= 1   
     \]
\end{proof}
\begin{exercise}{}{}
设随机变量 $X$ 的分布函数为
\[
F(x) = 1 - e^{-x^2}, \quad x>0,
\]
试求 $E(X)$ 和 $\operatorname{Var}(X)$.
\end{exercise}
\begin{proof}
   \(  F\left( 0 \right)= \lim_{x\to 0^{+ }}F\left( x \right)= 0    \), 因此 \(  P\left( X\le 0 \right)= 0   \), 从而 \(  E\left( X^{-} \right)= 0   \).   由Exercise \ref{ex:10-21-1} (1) \[
    \begin{aligned}
    E\left( X \right)=   E\left( X ^{+ }\right)&= \int_{0}^{\infty}P\left( X^{+ }> x \right) \,\mathrm{d} x\\ 
      &= \int_{0}^{\infty}P\left( X> x \right)\,\mathrm{d} x\\  
       &= \int_{0}^{\infty}\left( 1-F\left( x \right)  \right)\,\mathrm{d} x\\ 
        &= \int_{0}^{\infty}e^{-x^{2}}\,\mathrm{d} x\\ 
         &= \frac{\sqrt{ \pi } }{2 }   
    \end{aligned}
     \] 由Exercise \ref{ex:10-21-1}  (2) \[
     \begin{aligned}
     E\left( X^{2} \right)= E\left( \left( X^{+ } \right)^{2}  \right)&=    \int_{0}^{\infty}2xP\left( X^{+ }> x \right)\,\mathrm{d} x \\ 
      &= \int_{0}^{\infty}2xP\left( X> x \right)\,\mathrm{d} x\\ 
       &= \int_{0}^{\infty}2x\left( 1-F\left( x \right)  \right)\,\mathrm{d} x\\ 
        &= \int_{0}^{\infty}2xe^{-x^{2}}\,\mathrm{d} x\\ 
         &= \int_{0}^{\infty}e^{-x^{2}}\,\mathrm{d} x^{2}\\ 
          &= 1  
     \end{aligned}
     \]于是 \[
     \mathrm{Var}\left( X \right)=  E\left( X^{2} \right)-\left( E\left( X \right)  \right)^{2}= 1- \frac{ \pi  }{4 }    
     \]
\end{proof}
\begin{exercise}{}{10-21-2}
试证:对任意的常数 $c \neq E(X)$, 有
\[
\operatorname{Var}(X) =E[\left( X-E\left( X \right)  \right)^{2} ]< E[\left( X-c \right)^{2}  ]
\]
\end{exercise}
 \begin{proof}
   \[
   \begin{aligned}
   E\left[ \left( X-c \right)^{2}  \right]-E\left[ \left( X-E\left( X \right)  \right)^{2}  \right]&= E\left[ \left( X-c \right)^{2}- \left( X-E\left( X \right)  \right)^{2}   \right]\\ 
    &= E\left[ \left( E\left( X \right)-c  \right)\left( 2X-c-E\left( X \right)  \right)   \right]\\ 
     &= \left( E\left( X \right)-c  \right)E[\left( 2X-c-E\left( X \right)  \right) ]\\ 
      &= \left( E\left( X \right)-c  \right)\left( 2E\left( X \right)-c-E\left( X \right)   \right)\\ 
       &= \left( E\left( X \right)-c  \right)^{2}> 0     
   \end{aligned}
   \]
 \end{proof}
\begin{exercise}{}{}
设随机变量 $X$ 仅在区间 $[a,b]$ 上取值, 试证
\[
a \le E(X) \le b, \operatorname{Var}(X) \le \left(\frac{b-a}{2}\right)^2.
\]
\end{exercise}
\begin{proof}
    对于第一个不等式, 易见 \[
    a= \int a \,\mathrm{d} P\le \int X\,\mathrm{d} P= E\left( X \right)\le \int b\,\mathrm{d} P= b 
    \]
   由 Exercise \ref{ex:10-21-2}, 取 \(  c= \frac{a+ b }{2 }   \), 则   \[
      \mathrm{Var}\left( X \right)\le  E\left[ \left( X-\frac{a+ b }{2 }  \right)^{2}  \right]  
      \] 由于 \[
      \left| X-\frac{a+ b }{2 }  \right|\le \max \left\{ \left| b-\frac{a+ b }{2 }  \right|, \left| a-\frac{a+ b }{2 }  \right|   \right\} = \frac{b-a }{2 } 
      \]我们有 \[
      \left( X-\frac{a+ b }{2 }  \right)^{2}\le \left( \frac{b-a }{ 2}  \right)^{2}  
      \]从而 \[
      \mathrm{Var}\left( X \right)\le E\left[ \left( \frac{b-a }{ 2}  \right)^{2}  \right]= \left( \frac{b-a }{2 }  \right)^{2}   
      \]
\end{proof}

\begin{exercise}{}{10-23-1}
设 $g(x)$ 为随机变量 $X$ 取值的集合上的非负不减函数,且 $E(g(X))$ 存在,证明:对任意的 $\varepsilon>0$,有
\begin{equation}
P(X > \varepsilon) \leqslant \frac{E(g(X))}{g(\varepsilon)}.
\end{equation}
\end{exercise}
 \begin{proof}
    \[
    \begin{aligned}
    E\left( g\left( X \right)  \right)&= \int g\left( X \right)\,\mathrm{d} P\\ 
     &\ge \int_{ X>  \varepsilon } g\left( X \right)\,\mathrm{d} P    \\ 
      &\ge \int_{X>  \varepsilon }g\left(  \varepsilon  \right) \,\mathrm{d} P\\ 
       &= g\left(  \varepsilon  \right)\int_{X>  \varepsilon }\,\mathrm{d} P\\ 
        &= g\left(  \varepsilon  \right)P\left( X>  \varepsilon  \right)   
    \end{aligned}
    \]
 \end{proof}
\begin{exercise}{}{}
设 $X$ 为非负随机变量,$a>0$.若 $E(e^{aX})$ 存在,证明:对任意的 $x>0$,有
\begin{equation}
P(X \geqslant x) \leqslant \frac{E(e^{aX})}{e^{ax}}.
\end{equation}
\end{exercise}
\begin{proof}
由于 \(  e^{ay}  \)是不减的, 对于任意的 \(x> 0 \), \(  e^{ay}\ge e^{ax}, \forall y\ge x  \). 于是       \[
    \begin{aligned}
  E\left( e^{aX} \right)&= \int e^{aX}\,\mathrm{d} P \\ 
   &\ge \int_{X\ge x} e^{aX}\,\mathrm{d} P\\ 
    &\ge  \int_{X\ge x}e^{a x}\,\mathrm{d} P= e^{ax} \int_{X\ge x}\,\mathrm{d} P\\ 
     &= e^{ax} P\left( X\ge x \right) 
    \end{aligned}
    \]
\end{proof}
\begin{exercise}{}{}
已知正常成年男性每升血液中的白细胞数平均是 $7.3\times10^9$,标准差是 $0.7\times10^9$.试利用切比雪夫不等式估计每升血液中的白细胞数在 $5.2\times10^9$ 至 $9.4\times10^9$ 之间的概率的下界.
\end{exercise}
\begin{solution}
       由于 \[
    7.3\times 10^{9}= \frac{5.2\times 10^{9}+ 9.4\times 10^{9} }{2 } 
    \]取 \(  \mu = 7.3\times 10^{9}  \), \(   \sigma = 0.7\times 10^{9}  \), \(  X  \)为白细胞数,  由Chebyshev不等式, 所求概率满足   
    $$P(|X - \mu| \geq k\sigma) \leq \frac{1}{k^2}$$从而 \[
   P\left( \left| X-\mu  \right|< k  \sigma   \right)\ge 1-\frac{1 }{k^{2} }  
   \]这里 \[
   k= \frac{2.1\times 10^{9} }{ \sigma  }= \frac{2.1\times 10^{9} }{0.7\times 10^{9} }= 3  
   \], 得到下界为 \(  \frac{8 }{9 }   \)  
\end{solution}

\begin{exercise}{}{}
一批产品中有10\%的不合格品,现从中任取3件,求其中至多有一件不合格品的概率.
\end{exercise}
\begin{solution}
     设 \(  X  \)为取得的不合格品的件数, 则  \(  X\sim B\left( n,p \right)=   B\left( 3, 0.1 \right)   \) .至多有一件不合格品的概率为 \[
    P\left( X= 0 \right)+ P\left( X= 1 \right)=  0.9^{3}+ 3 \times 0.1\times 0.9^{2}= 0.972
    \]
\end{solution}

\begin{exercise}{}{}
一条自动化生产线上产品的一级品率为0.8,现检查5件,求至少有2件一级品的概率.
\end{exercise}
\begin{solution}
    设 \(  X  \)为一级品的件数, 则 \(  X\sim  B\left( 5, 0.8 \right)    \). 则 \[
   \begin{aligned}
    P\left( X<  2 \right)&=\binom{ 5 }{0  } 0.2^{5}+ \binom{ 5 }{1  }0.8\times 0.2^{4}=0.00672
   \end{aligned}
    \]因此至少有两件一级品的概率为 \[
    P\left( X\ge 2 \right)= 1-P\left( X< 2 \right)= 0.99328
    \]
\end{solution}

\begin{exercise}{}{}
设随机变量 $X \sim b(n,p)$, 已知 $E(X) = 2.4$, $\mathrm{Var}(X) = 1.44$, 求两个参数 $n$ 与 $p$ 各为多少?
\end{exercise}
\begin{proof}
     \[
    np= E\left( X \right)= 2.4 
     \] \[
     np\left( 1-p \right)= \mathrm{Var}\left( X \right)=  1.44 
     \]从而 \[
     \left( 1-p \right)= \frac{1.44 }{2.4 }= 0.6\implies p= 0.4  
     \]进而  \[
     n = \frac{2.4 }{p }= 6 
     \]综上 \[
     n = 6,\quad p= 0.4
     \]
\end{proof}
\begin{exercise}{}{}
设随机变量 $X$ 服从二项分布 $b(2,p)$, 随机变量 $Y$ 服从二项分布 $b(4,p)$. 若 $P(X \geqslant 1) = 8/9$, 试求 $P(Y \geqslant 1)$.
\end{exercise}
\begin{solution}
      \[
     \binom{ 2 }{0  }\left( 1-p \right)^{2}= P\left( X= 0 \right)=    P\left( X< 1 \right)= 1-P\left( X\ge  1 \right)= \frac{1 }{9 }   
     \]得到 \[
     \left( 1-p \right)^{2}= \frac{1 }{9 }  
     \]又\[
     \binom{ 4 }{0  } \left( 1-p \right)^{4}=P\left( Y= 0 \right)=   P\left( Y< 1 \right)= 1-P\left( Y\ge 1 \right)   
     \] 因此 \[
     P\left( Y\ge 1 \right)=1- \left( \left( 1-p \right)^{2}  \right)^{2}= 1- \left( \frac{1 }{9 }  \right)^{2}= \frac{80 }{81 }   
     \]
\end{solution}

\begin{exercise}{}{}
一批产品的不合格品率为 $0.02$, 现从中任取 $40$ 件进行检查, 若发现两件或两件以上不合格品就拒收这批产品. 分别用以下方法来拒收的概率:
\begin{enumerate}
\item 用二项分布作精确计算;
\item 用泊松分布作近似计算.
\end{enumerate}
\end{exercise}
\begin{solution}
      \begin{enumerate}
        \item 设 \(  X  \)为不合格品的数量, 则 \(  X\sim B\left( 40,0.02 \right)   \), \[
        P\left( X< 2 \right)= \binom{ 40 }{0  }0.02^{0}\times 0.98^{40}+ \binom{ 40 }{1  }0.02\times 0.98^{39}   
        \]  拒收的概率为 \[
        P\left( X\ge 2 \right)= 1-P\left( X< 2 \right)\approx  0.1904
        \]
        \item 由于$n=40$较大而$p=0.02$较小, \(  np= 0.8  \)  适中. 可以认为 \(  X  \)近似服从\(  \mathrm{Poisson}\left( 0.02 \times 40\right)= \mathrm{Poisson}\left( 0.8 \right)    \),   我们有 \[
        P\left( X< 2 \right)\approx  e^{-0.8}+ 0.8e^{-0.8} 
        \]于是 \[
        P\left( X\ge 2 \right)= 1-P\left( X< 2 \right)\approx   0.1912\]
    \end{enumerate}
    
\end{solution}

\begin{exercise}{}{}
设 $X$ 服从泊松分布, 且已知 $P(X=1) = P(X=2)$, 求 $P(X=4)$.
\end{exercise}
\begin{solution}
      \[
     P\left( X= k \right)= \frac{  \lambda ^{k}e^{- \lambda }}{k! }  
     \]由条件 \[
      \lambda e^{- \lambda }=  \frac{1 }{2 } \lambda ^{2}e^{- \lambda }
     \]得到 \[
      \lambda = 0,\quad \text{or},\quad  \lambda = 2
     \]由于泊松分布的定义 \(   \lambda > 0  \), 这里只取 \(   \lambda = 2  \), 从而 $$P(X=4) = \frac{\lambda^4 e^{-\lambda}}{4!}= \frac{16e^{-2}}{24} = \frac{2e^{-2}}{3}$$
\end{solution}

\begin{exercise}{}{}
已知某商场一天来的顾客数 $X$ 服从参数为 $\lambda$ 的泊松分布, 而每个来到商场的顾客购物的概率为 $p$, 证明: 此商场一天内购物的顾客数服从参数为 $\lambda p$ 的泊松分布.
\end{exercise}
\begin{proof}
   设 \(  Y  \)为一天内购物的顾客数. 则 \[
   P\left( Y= k \right)= \sum _{j= k}^{\infty}P\left( Y= k| X= j \right)P\left( X= j \right)    ,\quad k= 0,1,2,\cdots 
   \]  \(  Y|\left( X= j \right)   \)是一个服从 \(  B\left( j,p \right)   \)的随机变量, 我们有 \[
   P\left( Y= k|X= j \right)= \binom{ j }{k  }p^{k}\left( 1-p \right)^{j-k}   
   \]由于 \(  X  \)服从参数为 \(   \lambda   \)的泊松分布, 我们有 \[
   P\left( X= j \right)=  \frac{ \lambda ^{j} e^{- \lambda }}{j! } ,\quad j= 0,1,2,\cdots 
   \]  因此 \[
  \begin{aligned}
   P\left( Y= k \right)= \sum _{j= k}^{\infty}\binom{ j }{k  }p^{k}\left( 1-p \right)^{j-k} \frac{ \lambda ^{j} e^{- \lambda }}{j! }  &  = \sum _{j= k}^{\infty}\frac{ \lambda ^{j}e^{- \lambda }}{k!\left( j-k \right)!  }p^{k}\left( 1-p \right)^{j-k}   \\ 
    &= \sum _{i= 0}^{\infty}\frac{ \lambda ^{i+ k}e^{- \lambda } }{k! i! }p^{k}\left( 1-p \right)^{i}  \\ 
     &= \frac{ \lambda ^{k}p^{k}e^{- \lambda } }{k! }\sum _{i= 0}^{\infty}\frac{\left( \lambda \left( 1-p \right) \right)^{i}  }{i! }\\ 
      &= \frac{ \lambda ^{k}p^{k}e^{- \lambda } }{k! } e^{ \lambda \left( 1-p \right) }   \\ 
       &= \frac{\left(  \lambda p \right)^{k} e^{- \lambda p}  }{k! } \\ 
        &,\quad k= 0,1,2,\cdots 
  \end{aligned}
   \]即 \(  Y  \)服从 \(  \mathrm{Poisson}\left(  \lambda p\right)   \)  
\end{proof}
\begin{exercise}{}{}
设一个人一年内患感冒的次数服从参数 $\lambda = 5$ 的泊松分布. 现有某种预防感冒的药物对 $75\%$ 的人有效(能将泊松分布的参数减少为 $\lambda = 3$), 对另外的 $25\%$ 的人不起作用. 如果某人服用了此药, 一年内患了两次感冒, 那么该药对他(她)有效的可能性是多少?
\end{exercise}
\begin{solution}
    设该药有效的事件为 \(  A  \), 设一年两次感冒的事件为 \(  B  \), 则药有效的可能性为 \[
  P\left( A|B \right) = \frac{P\left( B|A \right)P\left( A \right)   }{ P\left( B \right) }=  \frac{P\left( B|A \right)P\left( A \right)   }{P\left( B|A \right)P\left( A \right)+ P\left( B|A^{c} \right)P\left( A^{c} \right)     } 
    \] 设 \(  X  \)为感冒的次数. 则  \(  X|A  \)服从 \(  \mathrm{Poisson}\left( 3 \right)   \), \(  X|A^{c}  \)服从 \(  \mathrm{Poisson}\left( 5 \right)   \). 我们有 \[
    P\left( B|A \right)= P\left( X|A= 2 \right)= \frac{3^{2}e^{-3} }{2! }= \frac{9 }{2 }e^{-3}    
    \] \[
    P\left( B|A^{c} \right)= P\left( X|A^{c}= 2 \right)= \frac{5^{2}e^{-5} }{2! }= \frac{5^{2} }{2}e^{-5}    
    \]     此外,  \[
    P\left( A \right)= \frac{3 }{4 },\quad P\left( A^{c} \right)= \frac{1 }{4 }    
    \]全部带入后计算得到 \[
    P\left( A|B \right)= \frac{\frac{27 }{8 }e^{-3}  }{\frac{27 }{8 }e^{-3}+ \frac{5 }{2 }e^{-5}\frac{1 }{4 }    }  =\frac{27e^{2} }{27e^{2}+ 25 } 
    \]
\end{solution}

\begin{exercise}{}{}
从一个装有 $m$ 个白球、$n$ 个黑球的袋中有返回地摸球，直到摸到白球时停止。试求取出黑球数的期望。
\end{exercise}

\begin{solution}
        设 \(  X  \)为取出黑球的个数. 则 \(  Y= X+ 1  \)  为第一次摸到白球时, 摸出的球的个数. 单次摸出白球的概率为 \(  p= \frac{m }{m+ n }   \). 因此 \(  Y  \)服从参数为 \(  \frac{m }{m+ n }   \)的几何分布. 我们有 \[
    E\left( Y \right)= \frac{1 }{p }= \frac{m+ n }{m }   
    \]由期望的线性, 取出黑球数的期望为 \[
    E\left( X \right)= E\left( Y-1 \right)= E\left( Y \right)-1=  \frac{n }{m }    
    \]   
\end{solution}

\begin{exercise}{}{}
某种产品上的缺陷数 $X$ 服从下列分布列：
\[
P(X = k) = \frac{1}{2^{k+1}}, \quad k = 0, 1, \cdots,
\]
求此种产品上的平均缺陷数。
\end{exercise}
\begin{solution}
   令 \(  Y= X+ 1  \), 则 \[
   P\left( Y= k \right)= P\left( X = k-1\right)=\frac{1 }{2^{k} }=   \left( \frac{1 }{2 } \right)^{k-1} \left( 1-\frac{1 }{2 }  \right)    ,\quad  k= 1,2,\cdots 
   \]       因此 \(  Y  \)服从参数为 \( p=  \frac{1 }{2 }   \)的几何分布. 我们有 \[
   E\left( Y \right)= \frac{1 }{p } = 2
   \]  由期望的线性,  \[
   E\left( X \right)= E\left( Y-1 \right)= 2-1= 1  
   \]为平均缺陷数.
\end{solution}

\begin{exercise}{}{}
某产品的不合格品率为 0.1,每次随机抽取 10 件进行检验,若发现其中不合格品数多于 1,就去调整设备.若检验员每天检验 4 次,试问每天平均要调整几次设备?
\end{exercise}
\begin{solution}
          设 \(  Y  \)为四次检验中, 不合格品多于1 的次数. 设 \(  X  \)为一次检验中, 不合格品的数量. 则  \(  X\sim B\left( 10, 0.1 \right)   \). 我们有 \[
    P\left( X\le 1 \right)= \binom{ 10 }{0  }0.1^{0}\times 0.9^{10}+ \binom{ 10 }{1  }0.1\times 0.9^{9}=  0.9^{10}+ 0.9^{9}  
    \]   从而 \[
    P\left( X> 1 \right)= 1-0.9^{10}-0.9^{9} 
    \]我们将上述值记作 \(  p  \), 为不合格品数多于1的概率. 从而 \(  Y\sim B\left( 4,p \right)   \). 平均每天调整设备的次数为 \[
    E\left( Y \right)=  4p= 4\left( 1-0.9^{10}-0.9^{9} \right) 
    \]  
\end{solution}

\begin{exercise}{}{}
一个系统由多个元件组成,各个元件是否正常工作是相互独立的,且各个元件正常工作的概率为 $p$.若在系统中至少有一半的元件正常工作,那么整个系统就有效.问 $p$ 取何值时,5 个元件的系统比 3 个元件的系统更有可能有效?
\end{exercise}
\begin{solution}
      设 \(  X  \)为5个原件中正常工作的元件个数, \(  Y  \)为 \(  3  \)个原件中正常工作的元件个数. 则 \(  X\sim B\left( 5,p \right)   \), \(  Y\sim B\left( 3,p \right)   \). \[
   \begin{aligned}
    P\left( X\ge 3 \right)&= \binom{ 5 }{3  }p^{3}\left( 1-p \right)^{2}+ \binom{ 5 }{4  }p^{4}\left( 1-p \right)+ \binom{ 5 }{5  }p^{5} \\ 
     &= 10p^{3}\left( 1-p \right)^{2}+ 5p^{4}\left( 1-p \right)+ p^{5}  
   \end{aligned}      
    \]     \[
    P\left( Y\ge 2 \right)= \binom{ 3 }{2  }p^{2}\left( 1-p \right) + \binom{ 3 }{3  }p^{3}=    3p^{2}\left( 1-p \right)+ p^{3} 
    \], \(  P\left( X\ge 3 \right)   >  P\left( Y\ge 2 \right) \) , 当且仅当 \[
   \begin{aligned}
    &10p^{3}\left( 1-p \right)^{2}+ 5p^{4}\left( 1-p \right)+ p^{5}> 3p^{2}\left( 1-p \right)+ p^{3}   \\ 
     &\iff 10p^{3}\left( 1-p \right)^{2}+ 5p^{4}\left( 1-p \right)>  3p^{2}\left( 1-p \right)+ p^{3}\left( 1-p \right)\left( 1+ p \right)    \\ 
      &\iff 10p^{3}\left( 1-p \right)+ 5p^{4}>  3p^{2}+ p^{3}\left( 1+ p \right)\\ 
       &\iff 10p^{3}-10p^{4}+ 5p^{4}>  3p^{2}+ p^{3}+ p^{4}\\ 
          &\iff 2p^{2}-3p+ 1<  0\\ 
           &\iff \left( 2p-1 \right)\left( p-1 \right) <  0\\ 
            &\iff \frac{1 }{2 }< p< 1 
   \end{aligned}
    \] 因此当 \(  \frac{1}{2}< p< 1  \)时, 更有可能有效. 
\end{solution}

\begin{exercise}{}{}
设随机变量 $X$ 服从参数为 $\lambda$ 的泊松分布,试证明
\[
E(X^n) = \lambda E[(X+1)^{n-1}],
\]
利用此结果计算 $E(X^3)$.
\end{exercise}
\begin{solution}
   \[
   \begin{aligned}
   E\left( X^{n} \right)&= \sum _{k= 0}^{\infty}k^{n}P\left( X= k \right)=  \sum _{k= 0}^{\infty} k^{n}\frac{ \lambda ^{k}e^{- \lambda } }{k! }  \\ 
    &= \sum _{k= 1}^{\infty}k^{n}\frac{ \lambda ^{k}e^{- \lambda } }{k! }\\ 
     &= \sum _{ k= 0}^{\infty} \left( k+ 1 \right)^{n}\frac{ \lambda ^{k+ 1}e^{- \lambda } }{\left( k+ 1 \right)!  }   \\ 
      &=  \lambda \sum _{k= 0}^{\infty}\left( k+ 1 \right)^{n-1}\frac{ \lambda ^{k}e^{- \lambda } }{k! }  \\ 
       &=  \lambda \sum _{k= 0}^{\infty}\left( k+ 1 \right)^{n-1}P\left( X= k \right)\\ 
        &=  \lambda E\left( \left( X+ 1 \right)^{n-1}  \right)    
   \end{aligned}
   \] \[
  \begin{aligned}
   E\left( X^{3} \right)= \lambda E\left( \left( X+ 1 \right)^{2}  \right)&=  \lambda  E\left( X^{2} \right)+ 2 \lambda E\left( X \right)+  \lambda  \\ 
    &=  \lambda ^{2}E\left( \left( X+ 1 \right)  \right)+ 2 \lambda E\left( X \right)+  \lambda \\ 
     &= \left(  \lambda ^{2}+ 2 \lambda  \right)E\left( X \right)+  \lambda ^{2}+  \lambda \\ 
      &= \left(  \lambda ^{2}+ 2 \lambda  \right) \lambda +  \lambda ^{2}+  \lambda \\ 
       &=  \lambda ^{3}+ 3 \lambda ^{2}+  \lambda      
  \end{aligned}   
   \]
\end{solution}

\begin{exercise}{}{}
令 $X(n,p)$ 表示服从二项分布 $b(n,p)$ 的随机变量,试证明:
\[
P(X(n,p) \leq i) = 1 - P(X(n,1-p) \leq n-i-1).
\]
\end{exercise}
\begin{proof}
      \[
  P\left( X\left( n,p \right)\le i  \right)= \sum _{ k= 0}^{i}\binom{ n }{k }p^{k}\left( 1-p \right)^{n-k}   
  \] \[
  P\left( X\left( n,1-p \right)\le n-i-1  \right)= \sum _{ j= 0}^{n-i-1}\binom{ n}{j }\left( 1-p \right)^{j}p ^{n-j} 
  \]做指标替换 \(  k=n-j \) , 则上式化为 \[
  \sum _{k= i+ 1}^{n}\binom{ n }{n-k  } \left( 1-p \right)^{n-k}p^{k} = \sum _{k= i+ 1}^{n}\binom{ n }{k  }p^{k}\left( 1-p \right)^{n-k}  
  \]因此 \[
  \begin{aligned}
  &P\left( X\left( n,p \right)\le i  \right)+ P\left( X\left( n,1-p \right)\le n-i-1  \right)   \\ 
   &=  \sum _{k = 0}^{i}\binom{ n }{k  }p^{k}\left( 1-p \right)^{n-k}  + \sum _{k= i+ 1}^{n}\binom{ n }{k  }p^{k}\left( 1-p \right)^{n-k}\\ 
    &= \sum _{k = 0}^{n}\binom{ n }{k  }p^{k}\left( 1-p \right)^{n-k}\\
    &= \left( p+ \left( 1-p \right)  \right)^{n}\\ 
     &= 1     
  \end{aligned}
  \]
\end{proof}
\begin{exercise}{}{}
设随机变量 $X$ 服从参数为 $p$ 的几何分布,试证明:
\[
E\left(\frac{1}{X}\right) = \frac{-p\ln p}{1-p}.
\]
\end{exercise}
\begin{proof}
     由于 \(  0< \left( 1-p \right)< 1   \),  由$\ln(1-x)$的Taylor展开,  \[
      \ln \left( p \right)=\ln \left( 1-\left( 1-p \right)  \right)=   -\sum _{k = 1}^{\infty} \frac{\left( 1-p \right)^{k}  }{k} 
      \]因此 \[
     \begin{aligned}
      E\left( \frac{1 }{X }  \right)&= \sum _{k = 1}^{\infty}\frac{\left( 1-p \right)^{k-1}p  }{k }\\ 
       &= \frac{-p }{1-p }\left( - \sum _{k= 1}^{\infty}\frac{\left( 1-p \right)^{k}  }{k }  \right)\\ 
        &= \frac{-p\ln p }{1-p }     
     \end{aligned} 
      \]
\end{proof}
\begin{exercise}{}{}
设随机变量 $X \sim b(n,p)$,试证明
\[
E\left(\frac{1}{X+1}\right) = \frac{1-(1-p)^{n+1}}{(n+1)p}.
\]
\end{exercise}
\begin{proof}
    \[
    E\left( \frac{1 }{X+ 1 }  \right)=  \sum _{k = 0}^{n}\frac{1 }{k+ 1 }\binom{ n }{k  }   p^{k}\left( 1-p \right)^{n-k} 
    \]注意到 \[
    \frac{1 }{k+ 1 }\binom{ n }{k  }= \frac{1 }{k+ 1 }\frac{n! }{k!\left( n-k \right)!  }= \frac{1 }{n+ 1 }\frac{\left( n+ 1 \right)!  }{\left( k+ 1 \right)! \left( \left( n+ 1 \right)-\left( k+ 1 \right)   \right)!   }= \frac{1 }{n+ 1 }\binom{ n+ 1 }{k+ 1  }        
    \]因此 \[
   \begin{aligned}
    E\left( \frac{1 }{X+ 1 }  \right)&= \sum _{k = 0}^{n}\frac{1 }{n+ 1 }\binom{ n+ 1 }{k+ 1  }p^{k}\left( 1-p \right)^{n-k}    \\ 
     &= \frac{1 }{n+ 1 } \frac{1}{p } \sum _{l = 1}^{n+ 1} \binom{ n+ 1 }{l  }p^{l} \left( 1-p \right)^{n+ 1-l} \\ 
      &= \frac{1 }{n+ 1 }\frac{1 }{p }\left( 1- \left( 1-p \right)^{n+ 1}  \right) 
   \end{aligned}
    \]
\end{proof}

(10分) 在X是连续随机变量的条件下, 证明 $\operatorname{E}(X-a)_+ = \int_a^\infty P(X > x) dx$, 其中 $a$ 是常数,
$(x-a)_+ = \begin{cases} x-a, & x \ge a, \\ 0, & x < a. \end{cases}$
\begin{proof}
$(x-a)_+ = 1_{\{x>a\}}(x-a)$, 称为\dfntxt{非负函数}.
$\operatorname{E}(X-a)_+ = \int 1_{\{x>a\}}(x-a) dP_X$
.
由定义有
$P(X>x) = \int 1_{\{X: X>x\}} dP_X$
其中 $1$ 为示性函数. $P_X$ 为 $X$ 的推出测度.
对上式两侧从 $a$ 到 $\infty$ 积分, 得
$\int_a^\infty P(X>x) dx = \int_a^\infty \int 1_{\{X>x\}} dP_X dx$
由 Fubini 定理, 右侧可等于
$\begin{aligned} & \quad \int \int_a^\infty 1_{\{X:X>x\}} dx dP_X \\ &= \int \left( \int_a^{\left( X-a \right)_{+ }+ a } 1 dx \right) dP_X \\ &= \int (X-a)_+ dP_X \\ &= \operatorname{E}((X-a)_+) \end{aligned}$
故 $\int_a^\infty P(X>x) dx = \operatorname{E}(X-a)_+$.
\end{proof}
\end{document}